% 99-64(99쪽) — 이웃하는 두 변이 3, 4 인 평행사변형 ABCD 와 두 대각선이 이루는 각 $120^\circ$. 스캔 화질이 낮아 TikZ 로 다시 그림.
% 라벨은 원문 것만: 꼭짓점 A·B·C·D, 길이 3·4, 각 $120^\circ$(원문처럼 지시선으로 가리킨다). 넓이(답)는 적지 않는다.
\begin{tikzpicture}[line width=0.6pt, x=1cm, y=1cm]
  \coordinate (A) at (1.23,0.85);
  \coordinate (B) at (0,0);
  \coordinate (C) at (2.15,-0.03);
  \coordinate (D) at (3.38,0.82);
  \coordinate (P) at (1.69,0.41);
  \fill[black!6] (A) -- (B) -- (C) -- (D) -- cycle;
  \draw (A) -- (B) -- (C) -- (D) -- cycle;
  \draw (A) -- (C);
  \draw (B) -- (D);
  % 길이 표시의 점선 호
  \draw[densely dashed, line width=0.35pt] (B) to[bend left=16] (A);
  \draw[densely dashed, line width=0.35pt] (A) to[bend left=14] (D);
  % 각 표시와 지시선
  \draw[line width=0.45pt] ($(P)+(13.6:0.26)$) arc (13.6:136.3:0.26);
  \draw[line width=0.35pt] (2.42,-0.19) -- (1.95,0.56);
  \node[font=\small, fill=white, inner sep=1pt] at (0.40,0.70) {$3$};
  \node[font=\small, fill=white, inner sep=1pt] at (2.31,1.14) {$4$};
  \node[font=\small, anchor=north west] at (2.44,-0.21) {$120^\circ$};
  \node[font=\small, anchor=south] at (1.22,0.93) {$\pt{A}$};
  \node[font=\small, anchor=east] at (-0.09,0.00) {$\pt{B}$};
  \node[font=\small, anchor=north] at (2.15,-0.11) {$\pt{C}$};
  \node[font=\small, anchor=west] at (3.46,0.84) {$\pt{D}$};
\end{tikzpicture}
